\makeabstract
    {
    3D Printing Metastructure Legs for High-Efficiency Thermoelectric Generators
    }
    {
    William Sober\textsuperscript{1}, Mingzhao Li\textsuperscript{2}, Ya Tang\textsuperscript{2}, Yan Li\textsuperscript{2}
    }
    {
    \textsuperscript{1} Amherst College, Amherst, MA\\
\textsuperscript{2} Thayer School of Engineering at Dartmouth College, Hanover, NH
    }
    {
    Global energy use is rising unsustainably, with excessive energy being lost as wasted heat in settings like data centers, automobile exhaust, and power plants. Thermoelectric generators (TEGs) make electricity from heat, and could recover this wasted energy on a large scale. However, current state-of-the art models suffer from high manufacturing costs and low efficiency. The goal of this project is to develop an optimal procedure and formulation to 3D print high-solid-loading (90 weight\%) Bi\textsubscript{0.5}Sb\textsubscript{1.5}Te\textsubscript{3} metastructure legs for use in a high-efficiency thermoelectric generator using direct ink writing (DIW). 
    }
    {
    Different combinations of organic binders were tested along with pre and post processing procedures to create the most printable slurry that maintained a high efficiency.
    }
    {
    Multiple slurries were created that could print through an improved nozzle size of .60 mm down from .84 mm. These slurries also achieved {\textasciitilde}90 weight\% loading. Sieving, ball milling, and planetary mixing proved to be effective at improving the printability of the slurries. Nozzle clogging remained the dominant problem for longer print times. A thermoelectric test yielded a figure of merit of .62 which is competitive with top of the field formulations. 
    }
    {
    This study demonstrates that direct ink writing (DIW) 3D printing has the potential to manufacture complex thermoelectric generators with high loading thermoelectric materials. The thermoelectric performance of these materials has been improved along with advances in printability. Their high efficiency provides a new direction for sustainable energy solutions.
    }

    \newpage
    
